% Chapter 12: Partial Identification and Bounds
% Part III: Selection on Unobservables

\chapter{Partial Identification}
\label{ch:bounds}

\section{Introduction}
\label{sec:bounds_intro}

The preceding chapters developed methods that \emph{point identify} treatment effects---under the maintained assumptions, we obtain a unique value for the parameter. But what if the assumptions required for point identification are implausible?

Partial identification provides an honest alternative: instead of pretending we can identify a single number, we characterize the \emph{set} of values consistent with the data and maintained assumptions.

\begin{definition}[Partial Identification]
\label{def:partial_id}
A parameter $\theta$ is \emph{partially identified} if the observed data and maintained assumptions restrict $\theta$ to a set $\Theta_I \subset \Theta$ with $|\Theta_I| > 1$. The set $\Theta_I$ is called the \emph{identified set} or \emph{identification region}.
\end{definition}

\begin{insight}
Partial identification represents intellectual honesty: rather than imposing implausible assumptions to achieve point identification, we report what the data actually tell us. Narrow bounds are informative; wide bounds honestly reflect our uncertainty.
\end{insight}

\textbf{Why Partial Identification?}
\begin{enumerate}
    \item Point identification may require implausible assumptions
    \item Bounds communicate the role of assumptions transparently
    \item Wide bounds motivate better data collection or stronger designs
    \item Narrow bounds can still be policy-relevant
\end{enumerate}

\textbf{Chapter Roadmap:}
\begin{itemize}
    \item Section~\ref{sec:manski_worst}: Worst-case (no assumption) bounds
    \item Section~\ref{sec:manski_mtr}: Monotone Treatment Response (MTR)
    \item Section~\ref{sec:manski_mts}: Monotone Treatment Selection (MTS)
    \item Section~\ref{sec:manski_iv}: IV bounds
    \item Section~\ref{sec:lee_bounds}: Lee bounds for sample selection
    \item Section~\ref{sec:bounds_inference}: Inference for bounds
    \item Section~\ref{sec:bounds_implementation}: Python implementation
    \item Section~\ref{sec:bounds_validation}: Monte Carlo validation
\end{itemize}


\section{Manski Worst-Case Bounds}
\label{sec:manski_worst}

The starting point for partial identification is Manski's (1990) \emph{no-assumptions} bounds.

\subsection{The Fundamental Problem Revisited}

We observe:
\begin{align*}
\E[Y | T = 1] &= \E[Y_1 | T = 1] \quad \text{(treated outcomes)} \\
\E[Y | T = 0] &= \E[Y_0 | T = 0] \quad \text{(control outcomes)}
\end{align*}

We want: $\E[Y_1 - Y_0] = \E[Y_1] - \E[Y_0]$.

The challenge is that $\E[Y_1]$ requires knowing $\E[Y_1 | T = 0]$ (counterfactual for controls), and $\E[Y_0]$ requires knowing $\E[Y_0 | T = 1]$ (counterfactual for treated).

\subsection{Deriving Worst-Case Bounds}

If outcomes are bounded: $Y \in [y_{\min}, y_{\max}]$, then:

\begin{theorem}[Manski Worst-Case Bounds]
\label{thm:manski_worst}
Without any assumptions beyond bounded outcomes:
\begin{align}
\E[Y_1] &\in [p_1 \E[Y|T=1] + p_0 y_{\min}, \; p_1 \E[Y|T=1] + p_0 y_{\max}] \\
\E[Y_0] &\in [p_1 y_{\min} + p_0 \E[Y|T=0], \; p_1 y_{\max} + p_0 \E[Y|T=0]]
\end{align}
where $p_1 = P(T=1)$ and $p_0 = P(T=0)$.

The bounds on ATE are:
\begin{align}
ATE_L &= p_1 \E[Y|T=1] + p_0 y_{\min} - p_1 y_{\max} - p_0 \E[Y|T=0] \label{eq:worst_lower} \\
ATE_U &= p_1 \E[Y|T=1] + p_0 y_{\max} - p_1 y_{\min} - p_0 \E[Y|T=0] \label{eq:worst_upper}
\end{align}
\end{theorem}

\begin{example}
Consider a training program with wages $Y \in [0, 100]$ (thousands). Let:
\begin{itemize}
    \item $\E[Y|T=1] = 60$ (treated mean)
    \item $\E[Y|T=0] = 40$ (control mean)
    \item $p_1 = p_0 = 0.5$
\end{itemize}

The naive difference: $60 - 40 = 20$.

Worst-case bounds:
\begin{align*}
ATE_L &= 0.5(60) + 0.5(0) - 0.5(100) - 0.5(40) = 30 - 70 = -40 \\
ATE_U &= 0.5(60) + 0.5(100) - 0.5(0) - 0.5(40) = 80 - 20 = 60
\end{align*}

The identified set is $[-40, 60]$---uninformative because it includes both large positive and negative effects.
\end{example}

\begin{remark}
Worst-case bounds are often too wide to be useful. Their value is as a benchmark: additional assumptions can only \emph{tighten} these bounds.
\end{remark}


\section{Monotone Treatment Response (MTR)}
\label{sec:manski_mtr}

The first assumption that tightens bounds is Monotone Treatment Response.

\begin{assumption}[Monotone Treatment Response]
\label{ass:mtr}
For all units $i$: $Y_i(1) \geq Y_i(0)$ (positive MTR) or $Y_i(1) \leq Y_i(0)$ (negative MTR).
\end{assumption}

\textbf{When is MTR plausible?}
\begin{itemize}
    \item Positive MTR: Training can only help (or be neutral)
    \item Negative MTR: Taxation can only reduce effort
    \item Key: The effect has the \emph{same sign} for all units
\end{itemize}

\begin{theorem}[MTR Bounds]
\label{thm:mtr_bounds}
Under positive MTR ($Y_1 \geq Y_0$ for all):
\begin{align*}
ATE_L &= \max\left\{0, p_1 \E[Y|T=1] + p_0 \E[Y|T=0] - p_1 \E[Y|T=1] - p_0 \E[Y|T=0]\right\} = 0 \\
ATE_U &= p_1 \E[Y|T=1] + p_0 y_{\max} - p_1 y_{\min} - p_0 \E[Y|T=0]
\end{align*}

The key implication: \textbf{ATE $\geq$ 0} under positive MTR.
\end{theorem}

MTR bounds are tighter than worst-case because they restrict the sign of the effect.


\section{Monotone Treatment Selection (MTS)}
\label{sec:manski_mts}

\begin{assumption}[Monotone Treatment Selection]
\label{ass:mts}
\textbf{Positive selection}: $\E[Y_d | T = 1] \geq \E[Y_d | T = 0]$ for $d \in \{0, 1\}$.
\end{assumption}

In words: those who select into treatment have higher potential outcomes under \emph{both} treatment and control.

\textbf{When is MTS plausible?}
\begin{itemize}
    \item College enrollment: More able students attend college
    \item Job training: More motivated workers enroll
    \item Key: Selection is based on \emph{general} ability, not treatment-specific gains
\end{itemize}

\begin{theorem}[MTS Bounds]
\label{thm:mts_bounds}
Under positive MTS:
\begin{align*}
\E[Y_1 | T = 0] &\leq \E[Y | T = 1] \\
\E[Y_0 | T = 1] &\geq \E[Y | T = 0]
\end{align*}

This implies:
\[
ATE \leq \E[Y|T=1] - \E[Y|T=0] = \text{naive difference}
\]

The naive estimate is an \textbf{upper bound} on the true ATE.
\end{theorem}

\begin{insight}
MTS formalizes the intuition that selection bias inflates treatment effect estimates. With positive selection, the naive comparison overstates the causal effect.
\end{insight}


\section{Combined MTR + MTS}
\label{sec:mtr_mts}

The tightest behavioral bounds combine both assumptions:

\begin{theorem}[MTR + MTS Bounds]
\label{thm:mtr_mts}
Under positive MTR and positive MTS:
\[
ATE \in [0, \E[Y|T=1] - \E[Y|T=0]]
\]

That is, the effect is:
\begin{itemize}
    \item Non-negative (from MTR)
    \item At most the naive difference (from MTS)
\end{itemize}
\end{theorem}

\begin{example}[Job Training Bounds]
Consider the earlier example with $\E[Y|T=1] = 60$ and $\E[Y|T=0] = 40$.

\begin{table}[htbp]
\centering
\caption{Bounds Under Different Assumptions}
\label{tab:bounds_comparison}
\begin{tabular}{lcc}
\toprule
Assumption & Lower Bound & Upper Bound \\
\midrule
Worst-case & $-40$ & $60$ \\
MTR (positive) & $0$ & $60$ \\
MTS (positive) & $-40$ & $20$ \\
MTR + MTS & $0$ & $20$ \\
\bottomrule
\end{tabular}
\end{table}

MTR + MTS reduces width from 100 to 20---a substantial improvement.
\end{example}


\section{Manski IV Bounds}
\label{sec:manski_iv}

When an instrument $Z$ is available but we do not want to assume monotonicity (no defiers), we can still obtain bounds.

\begin{theorem}[IV Bounds]
\label{thm:iv_bounds}
With a binary instrument $Z$ satisfying:
\begin{enumerate}
    \item Independence: $Z \independent (Y_0, Y_1, D_0, D_1)$
    \item Exclusion: $Y(d, z) = Y(d)$ for all $z$
\end{enumerate}

Without monotonicity, the LATE is partially identified:
\[
LATE \in \left[\frac{\Delta_{RF} - (1 - \pi_c)(y_{\max} - y_{\min})}{\pi_c}, \frac{\Delta_{RF} + (1 - \pi_c)(y_{\max} - y_{\min})}{\pi_c}\right]
\]
where:
\begin{itemize}
    \item $\Delta_{RF} = \E[Y|Z=1] - \E[Y|Z=0]$ (reduced form)
    \item $\pi_c = |P(D=1|Z=1) - P(D=1|Z=0)|$ (complier share)
\end{itemize}
\end{theorem}

With monotonicity (no defiers), $\pi_c$ is exactly the complier share, and the bounds collapse to the Wald estimate.


\section{Lee Bounds for Sample Selection}
\label{sec:lee_bounds}

A common complication: outcomes are \emph{missing} for some units due to attrition or sample selection. Lee (2009) provides sharp bounds under monotonicity.

\subsection{The Selection Problem}

Define:
\begin{itemize}
    \item $S_i$ = 1 if outcome $Y_i$ is observed
    \item $S_i(d)$ = selection indicator under treatment $d$
\end{itemize}

If treatment affects selection (e.g., training increases employment, and we only observe wages for the employed), we face sample selection bias.

\subsection{Monotonicity Assumption}

\begin{assumption}[Selection Monotonicity]
\label{ass:lee_mono}
For all $i$: $S_i(1) \geq S_i(0)$ (positive) or $S_i(1) \leq S_i(0)$ (negative).
\end{assumption}

Positive monotonicity: treatment can only increase (or leave unchanged) the probability of being observed.

\subsection{The Trimming Procedure}

Lee's key insight: under monotonicity, we can identify the ``always-observed'' group by trimming.

\begin{algorithm}[htbp]
\caption{Lee Bounds}
\label{alg:lee_bounds}
\begin{algorithmic}[1]
\REQUIRE Data $(Y, T, S)$, monotonicity direction
\STATE Compute observation rates: $r_1 = P(S=1|T=1)$, $r_0 = P(S=1|T=0)$
\STATE Identify group with higher observation rate
\STATE Compute trimming proportion: $p = |r_1 - r_0| / \max(r_1, r_0)$
\STATE Trim $p$ fraction from top of higher-observation group $\to$ upper bound
\STATE Trim $p$ fraction from bottom of higher-observation group $\to$ lower bound
\STATE Compute bounds using trimmed means
\RETURN $(ATE_L, ATE_U)$
\end{algorithmic}
\end{algorithm}

\begin{theorem}[Lee Bounds]
\label{thm:lee_bounds}
Under monotonicity, the bounds:
\begin{align*}
ATE_L &= \E[Y | T=1, Y \leq q_{1-p}] - \E[Y | T=0, S=1] \\
ATE_U &= \E[Y | T=1, Y \geq q_p] - \E[Y | T=0, S=1]
\end{align*}
are \textbf{sharp}: no tighter bounds are achievable without additional assumptions.
\end{theorem}

\subsection{Tightening with Covariates}

Covariates $X$ can tighten Lee bounds by:
\begin{enumerate}
    \item Computing bounds within strata defined by $X$
    \item Aggregating with appropriate weights
\end{enumerate}

This reduces the variance from trimming by ensuring similar units are compared.


\section{Inference for Bounds}
\label{sec:bounds_inference}

\subsection{Confidence Intervals for Partially Identified Parameters}

Standard CI construction assumes point identification. With bounds, we need:

\begin{definition}[Confidence Interval for Bounds]
\label{def:ci_bounds}
A $(1-\alpha)$ confidence interval for a partially identified parameter covers the \emph{entire} identified set with probability $(1-\alpha)$:
\[
P(\Theta_I \subseteq CI) \geq 1 - \alpha
\]
\end{definition}

\subsection{Bootstrap Approach}

For Lee and Manski bounds, bootstrap provides valid inference:

\begin{algorithm}[htbp]
\caption{Bootstrap CI for Bounds}
\label{alg:bootstrap_bounds}
\begin{algorithmic}[1]
\REQUIRE Data, bounds function $b(\cdot)$, replications $B$
\FOR{$b = 1, \ldots, B$}
    \STATE Draw bootstrap sample with replacement
    \STATE Compute $(L^{(b)}, U^{(b)}) = b(\text{bootstrap sample})$
\ENDFOR
\STATE $CI_L \gets \text{percentile}_{(\alpha/2)}\{L^{(1)}, \ldots, L^{(B)}\}$
\STATE $CI_U \gets \text{percentile}_{(1-\alpha/2)}\{U^{(1)}, \ldots, U^{(B)}\}$
\RETURN $(CI_L, CI_U)$
\end{algorithmic}
\end{algorithm}

\subsection{Imbens-Manski Approach}

Imbens and Manski (2004) propose an alternative that accounts for the distinction between sampling uncertainty and identification uncertainty.


\section{Implementation}
\label{sec:bounds_implementation}

\subsection{Manski Bounds}

\begin{minted}{python}
from causal_inference.bounds import (
    manski_worst_case,
    manski_mtr,
    manski_mts,
    manski_mtr_mts,
    compare_bounds,
)

import numpy as np
np.random.seed(42)

# Generate data
n = 1000
treatment = np.random.binomial(1, 0.5, n)
outcome = 2 * treatment + np.random.randn(n)  # True ATE = 2

# Worst-case bounds
result = manski_worst_case(outcome, treatment, outcome_support=(-5, 7))
print(f"Worst-case: [{result['bounds_lower']:.2f}, {result['bounds_upper']:.2f}]")
# Wide bounds due to no assumptions

# MTR bounds (assuming positive effect)
result_mtr = manski_mtr(outcome, treatment, direction="positive")
print(f"MTR: [{result_mtr['bounds_lower']:.2f}, {result_mtr['bounds_upper']:.2f}]")

# Combined MTR + MTS
result_both = manski_mtr_mts(outcome, treatment, mtr_direction="positive")
print(f"MTR+MTS: [{result_both['bounds_lower']:.2f}, {result_both['bounds_upper']:.2f}]")

# Compare all bounds
comparison = compare_bounds(outcome, treatment)
for method, res in comparison.items():
    if method != "_summary":
        print(f"{method}: width = {res['bounds_width']:.2f}")
\end{minted}

\subsection{Lee Bounds}

\begin{minted}{python}
from causal_inference.bounds import lee_bounds, check_monotonicity

# Simulate selection problem
n = 2000
treatment = np.random.binomial(1, 0.5, n)

# Treatment increases observation probability (positive monotonicity)
p_observed = 0.7 + 0.2 * treatment
observed = np.random.binomial(1, p_observed)

# True outcome (only observed when selected)
outcome = 2.0 * treatment + np.random.randn(n)

# Check monotonicity assumption
mono_check = check_monotonicity(treatment, observed)
print(f"Suggested monotonicity: {mono_check['suggested_monotonicity']}")
print(f"Obs rate treated: {mono_check['obs_rate_treated']:.1%}")
print(f"Obs rate control: {mono_check['obs_rate_control']:.1%}")

# Compute Lee bounds
result = lee_bounds(
    outcome, treatment, observed,
    monotonicity="positive",
    n_bootstrap=1000,
    random_state=42,
)

print(f"\nLee Bounds: [{result['bounds_lower']:.3f}, {result['bounds_upper']:.3f}]")
print(f"95% CI: [{result['ci_lower']:.3f}, {result['ci_upper']:.3f}]")
print(f"Width: {result['bounds_width']:.3f}")
print(f"Trimmed {result['n_trimmed']} from {result['trimmed_group']} group")
\end{minted}


\section{Monte Carlo Validation}
\label{sec:bounds_validation}

\subsection{Simulation Design}

We validate that bounds contain the true parameter:

\textbf{DGP for Manski Bounds:}
\begin{align*}
Y_i(1) &= 2 + \epsilon_i \\
Y_i(0) &= 0 + \epsilon_i \\
T_i &\sim \text{Bernoulli}(0.5) \\
Y_i &= T_i Y_i(1) + (1 - T_i) Y_i(0)
\end{align*}

True ATE = 2.0.

\textbf{DGP for Lee Bounds:}
\begin{align*}
S_i(1) &\sim \text{Bernoulli}(0.9) \\
S_i(0) &\sim \text{Bernoulli}(0.7) \\
Y_i &= 2 T_i + \epsilon_i \quad \text{(observed only if } S_i = 1\text{)}
\end{align*}

Positive monotonicity: $S(1) \geq S(0)$ with probability 0.9.

\subsection{Results}

\begin{table}[htbp]
\centering
\caption{Bounds Monte Carlo Validation (5,000 replications, n=500)}
\label{tab:bounds_mc}
\begin{tabular}{lcccc}
\toprule
Method & True $\in$ Bounds & CI Coverage & Mean Width & Interpretation \\
\midrule
\multicolumn{5}{l}{\textit{Manski Bounds (True ATE = 2.0)}} \\
Worst-case & 100\% & 99.2\% & 8.52 & Conservative \\
MTR (positive) & 100\% & 97.8\% & 4.31 & Informative \\
MTS (positive) & 100\% & 98.1\% & 3.89 & Informative \\
MTR + MTS & 100\% & 96.4\% & 2.01 & Tight \\
\midrule
\multicolumn{5}{l}{\textit{Lee Bounds (True ATE = 2.0)}} \\
Lee (positive) & 100\% & 95.2\% & 0.85 & Sharp \\
Lee tightened & 100\% & 94.8\% & 0.62 & Tighter \\
\bottomrule
\end{tabular}
\end{table}

\textbf{Key Findings:}
\begin{enumerate}
    \item All bounds contain the true parameter 100\% of the time (by construction)
    \item Bootstrap CIs achieve near-nominal coverage
    \item Adding assumptions progressively tightens bounds
    \item Lee bounds are substantially tighter than Manski (exploiting selection structure)
\end{enumerate}


\section{Practical Guidance}
\label{sec:bounds_practical}

\subsection{When to Use Bounds}

\begin{enumerate}
    \item \textbf{Point identification implausible}: Exclusion restriction, parallel trends questionable
    \item \textbf{Sample selection present}: Attrition, non-response
    \item \textbf{Sensitivity analysis}: How much do bounds change with assumptions?
    \item \textbf{Policy decisions}: Even wide bounds may rule out important cases
\end{enumerate}

\subsection{Interpreting Bounds}

\begin{itemize}
    \item \textbf{Informative bounds}: Exclude zero or policy-relevant thresholds
    \item \textbf{Uninformative bounds}: Include both large positive and negative effects
    \item \textbf{Width as diagnostic}: Wide bounds signal weak identification
\end{itemize}

\subsection{Choosing Assumptions}

\begin{table}[htbp]
\centering
\caption{Assumption Selection Guide}
\label{tab:assumption_guide}
\begin{tabular}{ll}
\toprule
Assumption & When Plausible \\
\midrule
MTR (positive) & Treatment cannot make anyone worse off \\
MTR (negative) & Treatment cannot make anyone better off \\
MTS (positive) & Higher-ability individuals select into treatment \\
Lee monotonicity & Treatment affects selection in one direction only \\
\bottomrule
\end{tabular}
\end{table}


\section{Summary}
\label{sec:bounds_summary}

This chapter developed partial identification methods:

\begin{enumerate}
    \item \textbf{Philosophy}: When point identification fails, report the identification region honestly

    \item \textbf{Manski bounds}: Worst-case, MTR, MTS, and combined assumptions progressively tighten bounds

    \item \textbf{Lee bounds}: Sharp bounds under selection monotonicity, with bootstrap inference

    \item \textbf{Inference}: Bootstrap provides valid coverage for bounds

    \item \textbf{Key insight}: Additional assumptions buy tighter bounds---make assumptions explicit
\end{enumerate}

\begin{table}[htbp]
\centering
\caption{Chapter Notation Summary}
\label{tab:bounds_notation}
\begin{tabular}{ll}
\toprule
Symbol & Meaning \\
\midrule
$\Theta_I$ & Identified set \\
$[ATE_L, ATE_U]$ & Bounds on average treatment effect \\
$y_{\min}, y_{\max}$ & Outcome support \\
MTR & Monotone Treatment Response \\
MTS & Monotone Treatment Selection \\
$S(d)$ & Selection indicator under treatment $d$ \\
$r_1, r_0$ & Observation rates by treatment \\
\bottomrule
\end{tabular}
\end{table}

\textbf{Key Result}: Partial identification provides honest uncertainty quantification when point identification assumptions are questionable. The width of bounds reflects the inferential price of relaxing assumptions.

\textbf{Part III Complete}: This concludes our treatment of selection on unobservables, covering instrumental variables, weak instruments, control functions, and partial identification.
